%\documentclass[aps,	prd, onecolumn, a4paper, 10pt, preprintnumbers,  tightenlines, 
	%notitlepage,floatfix, nofootinbib,longbibliographyprd,superscriptaddress]{revtex4-2}
\documentclass[preprint,aps,nofootinbib,a4paper,superscriptaddress,longbibliography,amsfonts, amssymb,amsmath,titlepage]{revtex4-1}
\ProvidesPackage{jheppub}[2013/03/21 r534]
\numberwithin{equation}{section} % JHEP style eqn 
%\input{CQGformat.input}
%s\documentclass[cqg]{iopart}
\usepackage{amssymb}
%\pdfoutput=1
%\usepackage{iopams, setstack}
\usepackage{bm}
\usepackage{bbm}
\usepackage{color}
\usepackage{microtype}
\usepackage{IEEEtrantools}
\usepackage{comment}
\usepackage{mathrsfs}
\usepackage{amsmath}
\usepackage{amsthm}
\usepackage{enumitem}
\usepackage[normalem]{ulem}
%\usepackage[numbers,square,sort&compress]{natbib}
\usepackage[pdftex]{hyperref}
%\bibliographystyle{iopart-num}
\usepackage[T1]{fontenc}

\newcommand{\PZ}[1]{\textcolor{red}{\textsf{[PZ: #1]}}}
\newcommand{\newtext}[1]{\textcolor{cyan}{#1}}

\def\half{\tfrac{1}{2}}
\def\pihalf{\tfrac{\pi}{2}}


\def\pheq{\phantom{=}}
\newcommand{\preindex}[2]{\,{}_{#2}#1}
\usepackage[toc,title,page,titletoc]{appendix}
%  vectors and diff Forms 
\def\bF{\boldsymbol{F}}
\def\bd{\boldsymbol{\mathrm{d}}}
\newcommand{\x}[1]{\xi_{#1}^{\circ}} % gauge vector

% GHP weight
\def\GHPwt{\ensuremath{\circeq}}
% NP SpinCoeffs
\def\kap{\kappa'}
\def\sip{\sigma'}
\def\rhp{\rho'}
\def\rhph{{\rho'}^\circ}
\def\tap{\tau'}
\def\tabp{\bar{\tau}'}
\def\bep{\beta'}
\def\epp{\epsilon'}
\def\rhp{\rho'}
\def\sib{\bar\sigma}
\def\epb{\bar\epsilon}
\def\kab{\bar\kappa}
\def\beb{\bar\beta}
\def\rhb{\bar\rho}
\def\pib{\bar\pi}
\def\alb{\bar\alpha}
\def\gab{\bar\gamma}
\def\mub{\bar\mu}
\def\lab{\bar\lambda}
\def\nub{\bar\nu}
\def\tab{\bar\tau}
% NP leg
\def\lmu{l^\mu}
\def\nmu{n^\mu}
\def\mmu{n^\mu}
\def\mbmu{{\bar m}^\mu}
\def\mb{\bar m}
%
% Held scalars
%
\newcommand{\Held}[1]{{#1}^{\circ}}
\def\tah{{\tau}^\circ}
\def\tabh{{\bar\tau}^\circ}
\def\tauH{\tau^{\circ}}
\def\taubH{\bar{\tau}^{\circ}}
\def\Psih{ {\Psi}^\circ}
\def\Psibh{ {\bar{\Psi}}^\circ}
\def\PsiH{\Psi^{\circ}}
\def\PsibH{\bar\Psi^{\circ}}
\def\rhoph{{\rho}^{\prime \circ} }
\def\rhopH{\rho^{\prime \circ}}
\def\rhopbh{{\bar\rho}^{\prime \circ} }
\def\tauHsq{\tau^{\circ}{}^{2}}
\def\Oh{{\Omega}^\circ}

% Nicer looking GHP Ops FROM GHZ PAPER 
\font\ec=ecrm0800 at 12pt
\def\thpH{\tilde{\th}'}
\def\edthH{\tilde\edth}
\def\edthpH{\tilde\edth'}
\def\th{\hbox{\ec\char'336}}
\def\edth{\hbox{\ec\char'360}}
\def\thp{\hbox{\ec\char'336}'}
\def\edthp{\hbox{\ec\char'360}'} 
\def\Dbar{\tilde \edth} % Held edth
\def\Nbar{\tilde \th'}
\def\Edth{\hat{\edth}} % raising operator
\def\EdthBar{\bar{\Edth}} % lowering operator on sYlm
%
%  SWS Harmonics
%
\newcommand{\C}[1]{\lambda_{l,#1}}
\newcommand{\thA}[1]{\theta^{A}_{#1}}
\def\thA{\ensuremath{\theta,\varphi}}
\def\Omp{\Omega_\p}
\newcommand{\CGCsqrt}[3]{\sqrt{\frac{(2{#1}+1)(2{#2}+1)(2{#3}+1)}{4\pi}}} 
\newcommand{\wthreej}[6]{%
  \ensuremath{%
  \begin{pmatrix}
      #1 & #2 & #3 \\
      #4 & #5 & #6
    \end{pmatrix}
  }%
}
% mode indices
\newcommand{\emm}{m}
\newcommand{\ess}{s}
\def\wlm{\ell \emm \omega}
% spin weighted spherical harmonics
\newcommand{\Y}[1]{{}_{#1}Y_{\ell, m}}
\newcommand{\bY}[1]{{}_{#1}\bar{Y}_{\ell, m}}
%\newcommand{\Ym}[2]{{}_{#1}Y_{l{#2}}}
\newcommand{\Ym}[2]{\, {}_{#1}Y_{\ell, #2}}
\newcommand{\Ylm}[3]{\, {}_{#1}Y_{#2, #3}}
\newcommand{\bYlm}[3]{\, {}_{#1} \bar{Y}_{#2, #3}}

% Sph. harmonic defs (Peter) 
\def\l{\ensuremath{\ell}}
\def\lm{{\ensuremath{\ell \emm}}}
% at the particle
\newcommand{\Yp}[1]{\, {}_{#1}Y^\textrm{p}_{\ell ,\emm}}
\newcommand{\Yplm}[3]{\, {}_{#1}Y^\textrm{p}_{#2 ,#3}}
\newcommand{\bYp}[1]{\, {}_{#1}\bar{Y}^\textrm{p}_{\ell, \emm}}
\def\YpL{ Y_{\ell}^{\p}}
\def\bYpL{\bar{Y}_{\ell}^{\p}}
\newcommand{\YpsL}[1]{{}_{#1} Y_{\ell}^{\p}}
\newcommand{\bYsL}[1]{ {}_{#1}\bar Y_{\ell}^\p}
%\newcommand{\C}[1]{\lambda_{\ell, #1}}
\newcommand{\Ltwo}[2]{\langle {#1} \vert {#2} \rangle_{L^2} }
\def\iStwo{\int_{\mathbb S^2}}
\def\dOm{\dd \Omega}

\newcommand{\xlm}[1]{\langle\xi_{#1}^{\circ}\rangle_{\ell \emm}} % gauge vector modes

% Helicity basis vectors
\def\hmp{{\hat{m}}^+}
\def\hmm{{\hat{m}}^-}
\def\nUL{{\hat n}^{\langle L \rangle}}
\def\uLL{{\hat n}_{\langle L \rangle}}
\def\hz{{\hat z}}

%
% Operators
%

% Chandrasekhar ops
\newcommand{\chandR}[2]{\,{}_{#1} \mathcal{L}^\dagger_{#2}{}}
\newcommand{\chandL}[2]{\,{}_{#1} \mathcal{L}^\dagger_{#2}{}}
% Teukolsky/Wald/Chrzanowski operators
\renewcommand{\S}{{\mathcal S}}
\newcommand{\E}{{\mathcal E}}
\newcommand{\T}{{\mathcal T}}
% Lie derivatives
\newcommand{\Lie}{{\mathcal L}}
\newcommand{\Lcsdag}[2]{ {}_{-{#1}}{\mathcal L}^{\dagger}_{#2}}
\newcommand{\Lcs}[2]{ {}_{{#1}}\mathcal L_{#2}}
%  derivative operators
\def\pd{\partial}
\def\dd{{\rm d}}
\newcommand{\dbd}[1]{\frac{\pd}{\pd #1}}
\renewcommand{\Re}{\operatorname{Re}}
\renewcommand{\Im}{\operatorname{Im}}
\def\pd{\partial}
\newcommand{\abs}[1]{\vert #1 \vert}
%
% Eqns.
%
\def\beq{\begin{equation}}
\def\eeq{\end{equation}}
\def\beqs{\begin{subequations}}
\def\eeqs{\end{subequations}}
\def\bal{\begin{align}}
\def\eal{\end{align}}
\def\non{\nonumber \\}
\def\beq{\begin{equation}}
\def\eeq{\end{equation}}
\def\ud{\mathrm{d}}
\def\beqs{\begin{subequations}}
\def\eeqs{\end{subequations}}


\def\half{\tfrac{1}{2}}

% indices 
\def\mn{\ensuremath{{\mu\nu}}}
\def\a{\alpha}
\def\b{\beta}
\newcommand{\stf}[1]{\langle {#1}\rangle}

% coordinates 
\def\vphi{\ensuremath{{\varphi}}}
\def\sz{\ensuremath{{\sqrt{1-z^2}}}}
\def\z{\zeta}
\def\zb{\bar\zeta}
\def\e{\eta}
\def\eb{\bar\eta}

% at the particle 
\def\p{{\rm p}}
\def\rhop{\rho_{\p}}
\def\rhobp{\bar{\rho}_{\p}}
\def\rp{r_{\p}}
\def\dOm{\ensuremath{\delta^2(\Omega-\Omega_\p)}}
\def\zp{z_\p}
\def\rhbp{\bar{\rho}_\p}

\def\rhom{\rho_{\rm min}}


\def\Dr{\Delta r}



% Spacing
\newcommand{\nls}{\phantom{\frac{1}{\rho}}}
\newcommand{\sforall}{\quad \forall}

\begin{document}

\title{GHZ Numeric Notes}


\author{Peter Zimmerman}
\email{zimmerator@protonmail.com}


\maketitle
\tableofcontents

\section{Introduction}\label{sec:introduction}


\subsection{Corrector Tensor}
The corrector tensor $x_{\mu\nu}$ is determined by the source $T_{ab}$ through
a hierarchy of equations descending from the $l^{a}$-sector (IRG variant) of Einstein's equations $l^{a}(\mathcal E_{ab}[x] = T_{ab}).$  In GHP form, these equations read \cite{Green:2019nam}
\beq
\label{eq:xmmb}
\begin{split}
 \rho^2 \th \left( \frac{\rhb}{\rho^3} \th \left[\frac{\rho}{\rhb} x_{m\mb} \right] \right)&= T_{ll}
\end{split}
\eeq
for $x_{m\mb}$, 
\beq
\label{eq:xnm}
\begin{split}
& \frac{\rho}{2(\rho+\rhb)} 
\th\left( (\rho+\rhb)^2\th \frac{x_{nm}}{\rho(\rho+\rhb)}\right)  \\
=&\ T_{lm} - \half\{(\th+\rho-\rhb)(\edth+\tab'-\tau) + 2\tab'(\th-2\rho) -
(\edth-\tau-\tab')\rhb +2\rho\tau\}x_{m\mb}.
\end{split}
\eeq
for $x_{nm}$, and
\beq
\label{eq:xnn}
\begin{split}
 &\half (\rho+\rhb)^2\th \left( \frac{1}{\rho+\rhb}x_{nn} \right)\\
&= T_{ln}- \Re\{ \edthp \edth  + (\tab'-\tau)\edthp + \left[\edth(\tab+\tau')\right] - 2 \tau'\tau - 2 \Psi_2 -2 \rhb\rho' + 2(\rho' \th + \rho\th') -  \th'\th \}  x_{m\mb} \\
&\ -\Re  \{\left[
(\th-2\rho)(\edthp-\tab) + (\tap+\tab)(\th+\rhb)
-2(\edthp-\tap)\rho-2\tab\th\right] x_{nm} \}
\end{split}
\eeq
for $x_{nn}$. 

In this note, we give explicit solutions to equations \eqref{eq:xmmb}, \eqref{eq:xnm}, and \eqref{eq:xnn} for a point particle \eqref{eq:stress pp} in terms of readily computable coefficients written in terms of the source.  

We work in mostly outgoing Kerr-Newman coordinates wherein equations \eqref{eq:xmmb}, \eqref{eq:xnm}, and \eqref{eq:xnn} reduce to a sequence of three ODEs along the outgoing null-generators $l^{\a} = \pd_{r}$ of increasing $r$ the affine parameter labeling cuts of $\mathcal I^{+}$, or in BL coordinates.  For initial data, we assume trivial conditions on $x_{\a\b}$ at $\mathcal{H}^-$, consistent with a bound orbit supported away from the horizon for all time.	

Throughout the note, we adopt the mostly minus metric ``west coast'' signature $(+,-,-,-)$ and geometrized units $G=c=1$. Greek letters denote spacetime  indices, lowercase Latin letters spatial indices, capital Latin letters denote quantities on the $\mathbb S^2$.

To avoid tracking factors of $\pi$, we  adopt the convention wherein the perturbed EFE reads
\begin{equation}
\mathcal{E}_{\mu\nu} = T_{\mu\nu}.
\end{equation}
Unless stated otherwise, all other conventions and notation will be chosen to match GHZ.

\section{Preliminaries}\label{sec:prelims}



\section{Calculation of $x_{\mu\nu}$}\label{sec:main calc}

\subsection{General solution in the source-free region}
The solution to \eqref{eq:xmmb} in the vacuum region  $r>r_{\rm max}$ is given by
\beq
x^>_{m\mb} = a^\circ \left( \frac{\rhb}{\rho}+\frac{\rho}{\rhb} \right) + b^\circ (\rho+\rhb), \qquad r>r_{\rm max}, \\
\eeq
where $\rho/\rhb+\rhb/\rho$ and $\rho+\rhb$  are homogeneous solutions. The integration constants $a^\circ$ and $b^\circ$ are determined by the source and will be given in a later section.

We next turn to the $nm$-component, as determined by \eqref{eq:xnm}. To integrate this equation, we write it first in Held form
\begin{align}\label{eq:D1}
 \frac{\rho}{2(\rho+\rhb)} 
\th\left( (\rho+\rhb)^2\th \frac{x_{nm}}{\rho(\rho+\rhb)}\right) = T_{lm} -\frac{1}{2} \left\{ \rhb
(\th+\rho-\rhb)\Dbar - \rhb(3\rhb+\rho)\tah\th+4\rho\rhb^2\tah\right\} x_{m\mb},
\end{align}
wherein we have used Eqs. \eqref{eq:Hops},\eqref{eq:th Dbar commute}, and \eqref{eq:Held coeffs}. 
Integrating the result give
\begin{align}\label{eq:x_nm sol}
x^>_{nm}= & \ \ \rhb\Oh\Dbar a^\circ+ \frac{(\rho+\rhb)^2}{\rho}\tah a^\circ
+\rhb \Dbar b^\circ + \rhb(\rho+\rhb)\tah b^\circ \nonumber \\
&\ +\frac{1-\Oh{}^2\rhb^2}{\rhb}c^\circ+\rho(\rho+\rhb) e^\circ, \qquad  r > r_{\rm max}.
\end{align}


We begin the calculation of $x_{nn}$ by recasting Eq.~\eqref{eq:xnn} in Held form
\begin{align}
\half (\rho+\rhb)^2\th \left( \frac{x_{nn}}{\rho+\rhb} \right)= T_{ln} + \Re\mathcal{U} x_{m\mb} + \Re[\mathcal{V} x_{nm}],
\end{align}
where 
\begin{align}\label{eq:Uxmmb}
&\mathcal{U}= 
 \rho\rhb\left[- \Dbar'\Dbar-\tabh (\rho+\rhb)\Dbar+\tabh\Dbar\th+\tah(\rhb-\rho)\Dbar' +\tah\Dbar'\th\right]  + \Nbar \th
 \non
&+\tfrac12 \left[ \Psih\left( \rho^2-\rhb^2\right)-2\rho\rhb(\rhb-\rho)\tah\tabh
-4\rhb \rhph(1+\rho\Oh\right)]\th 
- 2 \rho \Nbar \non
&+ \rho^2 \Psih (2\rho-\rhb) + 2\rhb^2\left(\rhph-\rho(2\rho-\rhb)\tah\tabh \right)- \rho\rhb(\rho-3\rhb)\rhph\Oh ,
\end{align}
and
\begin{align}\label{eq:Vxnm}
\mathcal{V}=-\rho\{ -3\rho\Dbar' + [\Dbar'-\tabh (2\rhb+\rho)]\th+2\rho^2\tabh\}
\end{align}

The solution is given by 
\begin{equation}
\label{xnn}
\begin{split}
x^>_{nn} =& 2 \Re \bigg\{ \half \frac{\rho + \bar{\rho}}{\rho^2} \Nbar a^\circ + (1 + \rho \bar{\rho} \Omega^{\circ 2})\Dbar' \Dbar a^\circ + 2 \left( \rho^2 + 2 \rho \bar{\rho} \right) \tau^\circ \bar{\tau}^\circ a^\circ\\
&+ \left( \frac{\rho^2}{\bar{\rho}} + \frac{3}{2} \rho \bar{\rho} \Omega^\circ \right) \Psih a^\circ - 2 \rho^{\prime \circ} a^\circ + \Nbar  b^\circ + 2 \rho \bar{\rho} \bar{\tau}^\circ \Dbar b^\circ + 2 \rho^2 \bar{\rho} \tau^\circ \bar{\tau}^\circ b^\circ + \rho^2 \Psih b^\circ\\
&+ \rho^2 \Dbar' e^\circ + 2 \rho^2 \bar{\rho} \bar{\tau}^\circ e^\circ - \frac{\bar{\rho}}{\rho} \left( 4 \Omega^\circ + \frac{1}{\rho} \right) \Dbar' c^\circ + 4 \frac{\bar{\rho}}{\rho} \bar{\tau}^\circ c^\circ + \left( \rho + \bar{\rho} \right) f^\circ \bigg\}
\end{split}
\end{equation}


\section{Point-particle source}\label{subsec:source}
Let us now take the perturbation to be sourced by a particle of mass $\mu$ on a worldline $\gamma$ parametrized  by proper time here denoted $\lambda$. We further denote the position and four velocity of the particle by $x_\p^{\mu}(\lambda)$ and 
$v^\mu=\dd x_\p^\mu/\dd \lambda$, respectively.
It carries the point particle stress energy distribution
\beq\label{eq:stress pp}
T_{\mu\nu}=\mu\int_\gamma v_\mu v_\nu \,\delta_4(x,x_\p(\lambda)) \dd \lambda.
\eeq
We consider a bound orbit with radii between $r_{\rm min} < r_\p  < r_{\rm max} $ and denote the coordinate location of the particle in  retarded coordinates as $x^\mu_\p = (u_\p, r_\p,z_\p, \varphi_\p)$.  

\subsection{Held Decomposition of the source}

With retarded KN time $-\infty < u_\p < \infty$ taken as the time parameter, we now rewrite the source as
\begin{align}\label{eq:stress pp Held}
T_{\mn} &=\int_{-\infty}^\infty \left[ \mu\frac{ v_\mu v_\nu}{\Sigma_\p v^u}  \, \delta^{2}_{\p}(\Omega)  \right] \delta_\p(u) \delta_\p(r)  \dd u_\p , \\
&= \Held T_{\mu\nu}(u,z,\varphi;\rp)\delta_\p(r), \\
&=  \Held J_{\mu\nu}(u;r_\p,z_\p,\varphi_\p)\delta_\p^2(\Omega)\delta_\p(r),
\end{align}
where in general $x^\p(u)$ after doing the $u_\p$ integration, and 
where (in the outgoing Kerr chart)
\begin{equation}\label{eq:J}
T^{\circ}_{\mn}  :=
\mu\frac{ v_\mu v_\nu}{\Sigma_\p \gamma_\p}  \, \delta^{2}_{\p}(\Omega), %=: j_{\mn}\delta(u-u_p) \, \delta(z-z_p)\delta(\vphi-\vphi_p)
\qquad \Held J_{\mu\nu} := \mu\frac{\gamma_\p \dot x_\mu \dot x_\nu}{\Sigma_\p },
\end{equation} 
with the coordinate four-velocity
\begin{equation}
    \dot x^\mu := \frac{\dd x^\mu}{\dd\lambda}\frac{\dd\lambda}{\dd u} = \frac{v^\mu}{\gamma_\p}
\end{equation}
and
$\delta^{2}_{\p}(\Omega)=\delta(z-\zp)\delta(\varphi-\varphi_{\p})$ and where 
$ v^\alpha \pd_\alpha u = \gamma_\p $. Here, $J_{\mu\nu}^\circ$ and $\Held T_{\mu \nu}$ are taken at $\rp(u)$, $\Sigma_\p = r_\p^2+a^2z_\p^2$ and it is understood that $r_\p =r_\p(\lambda(u_\p))$ etc.
We also introduce the quantities $T_\mu^\circ$ and $\Held J_\mu$ by
\begin{align}\label{eq:Tcirc hat}
T_{\mu \nu}l^{\nu} = T^{\circ}_\mu (u,z,\varphi)\delta_{\p}(r) = \Held J_\mu(u;z_\p,\varphi_\p)   \delta_{\p}(r) \delta_\p^2(\Omega)
\end{align}

\section{Spectral Solver}

\section{Windowing the puncture}


\end{document}